%!TEX root = ../template.tex
%%%%%%%%%%%%%%%%%%%%%%%%%%%%%%%%%%%%%%%%%%%%%%%%%%%%%%%%%%%%%%%%%%%
%% chapter1.tex
%% NOVA thesis document file
%%
%% Chapter with introduction
%%%%%%%%%%%%%%%%%%%%%%%%%%%%%%%%%%%%%%%%%%%%%%%%%%%%%%%%%%%%%%%%%%%

\chapter{Introduction}
\label{cha:introduction}

This first chapter serves as a brief introduction to the work initiated within
the scope of \gls{CDA} at \gls{GMV} and later continued directly under
\gls{GMV}. It outlines the context and motivation for the work carried out, the
problem being addressed, the proposed solution, and the objectives defined for
the internship and its immediate continuation.

\section{Context and Motivation}

\gls{GNSS} are a family of satellite-based positioning systems, usually
referred to as constellations, that provide geolocation and time information to
a receiver anywhere on Earth. There are several different \gls{GNSS}
constellations in operation today, including the United States' \gls{GPS},
Russia's \gls{GLONASS}, the European Union's \gls{Galileo}\cite{europaGalileo},
and China's \gls{BeiDou}. These systems consist of a network of satellites that
transmit signals to receivers, which then use those signals to determine their
location.

As reliance on \gls{GNSS} technology grows, so does the need for accurate and
reliable testing and validation of \gls{GNSS} receivers and the systems built
around them.\cite{europaWhatGNSS} A signal simulator is a tool used to
reproduce the \gls{GNSS} signals that would otherwise be transmitted by
navigation satellites and received by an
antenna.\cite{insidegnssVerificationTesting} Such simulators are essential for
testing and validating \gls{GNSS} receivers: they serve as a controlled and
repeatable signal source, enabling engineers to test receivers under a wide
range of scenarios without requiring outdoor tests or waiting for satellites to
reach a desired orbital position.\cite{spirentWhatGNSSGPS}

\gls{SDR} is a radio communication paradigm in which components traditionally
implemented in hardware are instead realised in software. This approach offers
greater flexibility and adaptability, as the same hardware platform can be
configured for different signal types without requiring physical
modifications.\cite{gnuradioRadio}\cite{wikipediaSoftwaredefinedRadio} In the
context of this internship, \gls{SDR} hardware provides the transmission
layer for \gls{GNSS} signal simulation: pre-computed \gls{I/Q} sample files,
produced by signal generation tools, are streamed to \gls{SDR} hardware for
transmission as \gls{RF} signals.

Controlling such hardware programmatically, however, requires interfacing with
low-level vendor APIs that expose hardware primitives directly, demanding
careful sequencing of calls, manual state management, and detailed knowledge of
hardware internals. Bridging this gap — between the high-level intent of an
engineer setting up a test scenario and the low-level reality of the hardware
API — is the central motivation for the work developed in this internship.

Additionally, \gls{RS232} is a standard for serial communication widely used in
industrial and laboratory instrumentation. In the context of this work, it
provides the physical communication channel between the developed system and
external equipment. Rather than carrying a bespoke application protocol, the
serial link is used to transport \gls{PPP}, thereby creating an IP channel over
which the higher-level control interface can operate.\cite{wikipediaRS232}

\section{Problem}

The \gls{DTAPI}, the C++ library provided by DekTec for interfacing with their
\gls{SDR} hardware, is a powerful but low-level interface. Configuring a device
requires precise sequencing of API calls, manual management of device state,
and direct handling of hardware-level concerns such as FIFO buffer sizing, DMA
configuration, and modulation parameter composition. Without a higher-level
abstraction, engineers are required to interact with these low-level details
directly, which significantly increases development effort, the risk of subtle
bugs, and the barrier to integrating \gls{SDR} hardware into automated test
workflows.

Furthermore, no existing open-source solution connects the output of \gls{GNSS}
signal generation tools to professional-grade \gls{SDR} hardware through a
clean, programmatic interface. This forces engineers into manual, error-prone
workflows that are difficult to automate and hard to integrate into larger
systems.

\section{Proposed Work}

The internship aimed to address this gap by developing a C++ abstraction layer
over the \gls{DTAPI}, exposed through a \gls{gRPC} interface, to enable
programmatic and remotely accessible control of \gls{SDR} hardware for
\gls{GNSS} signal transmission. A secondary objective was the development of an
\gls{RS232}-based communication path for integration with external equipment.
In the final architecture, this took the form of a \gls{PPP} link over the
serial channel, allowing the same gRPC interface to be reused over a standard
IP transport.

It should be noted that the scope of the internship evolved throughout its
duration, as described in \autoref{sec:scope-evolution}. The phases outlined
below reflect the original plan; the actual work carried out is documented in
\autoref{sec:workplan}.

\emph{Phase 1 --- \gls{SDR} Control:}

Development of a C++ API abstracting the \gls{DTAPI}, providing high-level
functions for device management, signal transmission, and multi-device
synchronisation, without requiring knowledge of hardware internals.

\emph{Phase 2 --- Parameter Management:}

Development of a parameter management system enabling users to configure and
manage \gls{SDR} transmission parameters efficiently, without resorting to
direct low-level API calls.

\emph{Phase 3 --- \gls{RS232} Connectivity:}

Design and implementation of an \gls{RS232}-based connectivity layer for
integration with external equipment, including serial interface configuration,
\gls{PPP} link establishment, and validation of the resulting IP channel.

\section{Objectives}

The main objectives of this internship were:

\begin{itemize}
      \item Acquire working knowledge of the target \gls{SDR} hardware architecture and its
            vendor API.
      \item Understand the \gls{RS232} communication standard and its application in
            instrument control, as well as the use of \gls{PPP} over serial
            links.
      \item Implement C++ modules for \gls{SDR} control via the \gls{DTAPI}, exposing a
            clean and programmatic interface through \gls{gRPC}.
      \item Develop a parameter management interface for test scenario configuration.
      \item Implement an \gls{RS232}-based communication path that allows
            external equipment to reach the control interface through a
            \gls{PPP}-backed IP channel.
      \item Validate the overall system functionality through functional and integration
            testing.
      \item Produce technical documentation and a final report covering the developed
            interfaces and modules.
\end{itemize}

\section{Work Plan \& Schedule}
\label{sec:workplan}

Table~\ref{tab:work_plan_schedule} outlines the originally proposed work plan
for the internship. Table~\ref{tab:actual_work_plan_schedule} reflects the
actual timeline, which deviated from the original plan due to changes in
project scope described in \autoref{sec:scope-evolution}. Note that between
each implementation phase there were continuous testing and code review periods
not explicitly represented in either table.

\begin{table}[htbp]
      \centering
      \caption{Proposed Work Plan \& Schedule for the Internship.}
      \label{tab:work_plan_schedule}
      \begin{tabular}{|p{3cm}|p{8cm}|c|c|}
            \hline
            \multicolumn{4}{|c|}{\textbf{Proposed Work Plan \& Schedule}}                                              \\
            \hline
            Phase                                                                & Task             & Duration & Dates \\
            \hline
            Onboarding and Studying                                              &
            Introduction to the signal generator architecture; study of \gls{SDR}
            APIs; \gls{RS232} protocol fundamentals; review of existing codebase
            and tooling at \gls{GMV}                                             &
            3 weeks                                                              & Mar 11 -- Mar 31                    \\
            \hline
            \gls{SDR} Control Modules                                            &
            Development of C++ API for \gls{SDR} control; implementation of signal
            transmission functions; integration with vendor API and unit testing &
            4 weeks                                                              & Apr 01 -- Apr 28                    \\
            \hline
            Parameter Management                                                 &
            Development of parameter management system; implementation of
            configuration interface; integration with \gls{SDR} control modules  &
            4 weeks                                                              & Apr 29 -- May 26                    \\
            \hline
            \gls{RS232} Connectivity                                             &
            Design and implementation of an \gls{RS232}-based connectivity path
            for external equipment, including \gls{PPP} link bring-up and IP
            communication validation                                             &
            3 weeks                                                              & May 27 -- Jun 16                    \\
            \hline
            Integration, Testing \& Documentation                                &
            Integration of all components; system testing and validation;
            preparation of technical documentation                               &
            4.5 weeks                                                            & Jun 17 -- Jul 17                    \\
            \hline
      \end{tabular}
\end{table}

\begin{table}[htbp]
      \centering
      \caption{Actual Work Plan \& Schedule for the Internship.}
      \label{tab:actual_work_plan_schedule}
      \begin{tabular}{|p{3cm}|p{8cm}|c|c|}
            \hline
            \multicolumn{4}{|c|}{\textbf{Actual Work Plan \& Schedule}}                                                                                                                                                                                                  \\
            \hline
            Phase                                                                                    & Task                                                                                                                                & Duration & Dates            \\
            \hline
            Onboarding and Studying                                                                  &
            Introduction to the signal generator architecture; study of \gls{SDR}
            APIs; \gls{RS232} protocol fundamentals; review of existing codebase
            and tooling at \gls{GMV}                                                                 &
            1 week                                                                                   & Mar 11 -- Mar 17                                                                                                                                                  \\
            \hline
            \gls{SDR} Control \& Parameter Management                                                &
            Development of C++ API for \gls{SDR} control; implementation of signal
            transmission functions; integration with vendor API; parameter
            management interface                                                                     &
            6.5 weeks                                                                                & Mar 18 -- May 2                                                                                                                                                   \\
            \hline
            \gls{RS232} Connectivity                                                                &
            Design and implementation of an \gls{RS232}-based connectivity path,
            including \gls{PPP} link establishment and gRPC communication over the
            resulting IP channel                                                                     &
            1 week                                                                                   & May 3 -- May 9                                                                                                                                                    \\
            \hline
            Integration with Proprietary Systems \& Documentation                                    &
            Integration of developed components with existing proprietary systems
            at \gls{GMV}; writing of technical documentation, installation scripts,
            and validation tooling                                                                   &
            6 weeks                                                                                  & May 10 -- Jun 21                                                                                                                                                  \\
            \hline
            Complete Testing Suite \& Final Report \& Further development of technical documentation & Create a complete testing suite to test component integration; finalize the report; continue development of technical documentation & 4 weeks  & Jun 22 -- Jul 17 \\
            \hline
      \end{tabular}
\end{table}

\section{Scope Evolution}
\label{sec:scope-evolution}

The scope of the internship evolved in three notable ways relative to the
original plan.

The first was a minor protocol change: the serial communication interface was
originally specified in the internship proposal as \gls{RS485}, but the
external equipment in use at \gls{GMV} communicated over \gls{RS232}. This
discrepancy arose from miscommunication between the involved parties and was
identified early in the implementation phase. As \gls{RS232} and \gls{RS485}
serve similar serial-communication roles at the physical layer, the impact on
the implementation was minimal, requiring only adjustments to the electrical
and interface configuration.

The second change concerned the communication architecture carried by the
serial link. An early implementation explored a custom application protocol on
top of \gls{RS232}; however, as integration requirements became clearer, that
approach was replaced by a \gls{PPP} link over the same physical medium. This
allowed the system to expose the existing gRPC control interface over a
standard IP channel, reducing duplicated protocol logic and aligning the serial
integration path with the rest of the software architecture.

The third change was more significant in terms of effort. As the developed
components matured, it became necessary to integrate them with existing
proprietary systems at \gls{GMV} that had not been part of the original scope.
This integration work was not planned at the outset and resulted in additional
unplanned effort, reflected in the actual work plan in
Table~\ref{tab:actual_work_plan_schedule}. The nature of these systems is
confidential and will not be described in detail in this report; their
existence is acknowledged where relevant to explain architectural decisions
made during the work.

\section{Continuation Beyond the Curricular Stage}
\label{sec:continuation-context}

This report is formally rooted in the curricular internship established between
\gls{GMV} and \gls{FCT}. However, because the initial objectives were reached
ahead of schedule, the work transitioned naturally into a continuation phase
carried out directly under \gls{GMV}, outside the university contract.

Some of the technical decisions described in this report, most notably the
migration to CMake, the replacement of the custom \gls{RS232} protocol with a
\gls{PPP}-based link, and the consolidation of gRPC over that resulting IP
channel, belong to that continuation period or were finalized there. They are
included because they superseded intermediate solutions developed during the
curricular stage and therefore better represent the final engineering outcome
of the project.

For that reason, this document distinguishes between contractual context and
technical outcome: the internship provides the academic framing and original
work plan, while the architecture described in the following chapters
corresponds to the consolidated solution reached across the internship and its
immediate continuation.
